% ============================================================
% SCHEDE PG — CASA DORIA (5 personaggi)
%
% @rpg-schema: <https://rpg-schema.org/instances/genova1507/crew-doria>
%   fitd:hasMember
%     <.../pg-ottaviano-doria> ,
%     <.../pg-fulvio-doria> ,
%     <.../pg-bianca-doria> ,
%     <.../pg-fra-matteo> ,
%     <.../pg-giacomo-moro> .
% ============================================================

\chapter{Schede PG — Casa Doria}
\index{schede PG!Doria}

% ─────────────────────────────────────────────────────────────

% ════════════════════════════════════════════════════════════
% PG 1 — OTTAVIANO DORIA
% @rpg-schema: <https://rpg-schema.org/instances/genova1507/pg-ottaviano-doria>
%   a fitd:Character ;
%   rdfs:label "Ottaviano Doria" ;
%   fitd:crew <.../crew-doria> ;
%   fitd:role "Il Capofamiglia / L'Ammiraglio" ;
%   fitd:action [ fitd:actionName "Hunt" ; fitd:dots 2 ] ;
%   fitd:action [ fitd:actionName "Study" ; fitd:dots 1 ] ;
%   fitd:action [ fitd:actionName "Survey" ; fitd:dots 2 ] ;
%   fitd:action [ fitd:actionName "Skirmish" ; fitd:dots 3 ] ;
%   fitd:action [ fitd:actionName "Wreck" ; fitd:dots 2 ] ;
%   fitd:action [ fitd:actionName "Command" ; fitd:dots 3 ] ;
%   fitd:action [ fitd:actionName "Consort" ; fitd:dots 2 ] ;
%   fitd:hasBond <.../pg-fulvio-doria> ;
%   fitd:hasBond <.../pg-bianca-doria> ;
%   fitd:hasBond <.../pg-fra-matteo> ;
%   fitd:hasBond <.../pg-giacomo-moro> .
% ════════════════════════════════════════════════════════════

\pgheader[DoriaBlue]{Ottaviano Doria}{Il Capofamiglia / L'Ammiraglio}{Doria}
\pgquote{Il mare non chiede permesso. Nemmeno io.}

\section*{Background \& Motivazione}

\begin{rulebox}{}
  \textbf{Background:} Cinquantadue anni, cicatrici da battaglie navali in acque greche
  quando aveva vent'anni. È il volto ufficiale della famiglia a Genova: quello che i
  francesi invitano ai banchetti e che sorride mentre pianifica la loro rovina.

  \medskip\noindent\textcolor{DoriaBlue}{\textbf{Motivazione:}} Vuole rivedere Genova
  libera prima di morire. Non per ideologia — per orgoglio. I Doria governano il mare
  da trecento anni. Non tollera che un re straniero dica loro dove attraccare.
\end{rulebox}

\section*{Attributi \& Azioni}


\begin{attributoblock}[DoriaBlue]{INSIGHT — Mente}
  \azione{Caccia}{Hunt}{2}{Individua bersagli, legge mappe militari}
  \azione{Studio}{Study}{1}{Analisi tattica, storia navale}
  \azione{Osserva}{Survey}{2}{Valuta situazioni, scansiona l'orizzonte}
  \azione{Ingegno}{Tinker}{0}{Meccanismi, artiglieria navale}
\end{attributoblock}


\begin{attributoblock}[DoriaBlue]{PROWESS — Corpo}
  \azione{Finezza}{Finesse}{1}{Precisione, duello con stile}
  \azione{Ombra}{Prowl}{0}{Infiltrazione, movimento non visto}
  \azione{Combattimento}{Skirmish}{3}{Scontro diretto, spada}
  \azione{Forza bruta}{Wreck}{2}{Sfondare, distruggere}
\end{attributoblock}


\begin{attributoblock}[DoriaBlue]{RESOLVE — Spirito}
      \azione{Percezione}{Attune}{0}{Leggere il non detto}
    \azione{Comando}{Command}{3}{Guidare uomini, imporre la propria volontà}
    \azione{Relazioni}{Consort}{2}{Frequentare i potenti, muoversi nei palazzi}
    \azione{Persuasione}{Sway}{1}{Convincere, manipolare, negoziare}
  \end{attributoblock}

\medskip
\begin{pgclockbox}{Stress \hfill \clock{9}{0}}
  \small\textit{Traumi: $\square$ Freddo \quad $\square$ Duro \quad $\square$ Ossessionato \quad $\square$ Paranoico \quad $\square$ Irrazionale \quad $\square$ Spezzato}
\end{pgclockbox}

\section*{Abilità Speciale}

% @rpg-schema: <.../pg-ottaviano-doria>
%   fitd:specialAbility [ rdfs:label "Autorità dell'Ammiraglio" ;
%     fitd:usesPerScore 1 ] .

\begin{doriabox}{Autorità dell'Ammiraglio}
  Quando Ottaviano impartisce un ordine in una situazione di crisi, tutti i PG
  del tavolo che lo eseguono tirano con \textbf{Posizione Controllata} invece
  di Rischiosa, indipendentemente dalle circostanze.

  \medskip\noindent\textit{Una volta per colpo. Deve essere un ordine esplicito pronunciato in scena.}
\end{doriabox}

\section*{Legami con i Compagni}

% @rpg-schema: fitd:hasBond annotations on each bond

\begin{table}[h]
\centering\renewcommand{\arraystretch}{1.4}
\begin{tabularx}{\linewidth}{@{}L{2.4cm} X@{}}
  \toprule\rowcolor{DoriaBlue}
  \textcolor{white}{\textbf{PG}} & \textcolor{white}{\textbf{Legame}} \\
  \midrule
  Ser Fulvio & Il mio luogotenente. Lo conosco da trent'anni. Sa mentire bene — e lo fa troppo spesso. \\
  \rowcolor{DoriaBlue!8}
  Bianca & Mia nipote. Più intelligente di chiunque. La proteggo — ma è lei che protegge me. \\
  Fra' Matteo & Il cappellano. Lo ho salvato dai turchi. Mi deve qualcosa che non potrà mai restituire. \\
  \rowcolor{DoriaBlue!8}
  Giacomo & Il contrabbandiere. Non mi fido di lui. È il migliore nel suo mestiere. \\
  \bottomrule
\end{tabularx}
\end{table}

\section*{Equipaggiamento}
\begin{goldlist}
  \item Spada da lato di fattura spagnola — un dono che non ricorda da chi
  \item Carte nautiche segrete del porto (con fondali reali, non quelli mostrati ai francesi)
  \item Sigillo ufficiale della famiglia — vale più di cento spade
  \item Fiaschetta d'argento con grappa delle Cinque Terre. Non offre mai da bere.
\end{goldlist}

\clearpage

\begin{secretbox}
  \textbf{Segreto:} Tre anni fa, durante una crisi di liquidità, Ottaviano ha
  ipotecato in segreto due delle quattro galee familiari con un banchiere
  genovese di simpatie francesi. Se la famiglia fallisce nel colpo riservato
  e non recupera denaro, il banchiere può confiscare le navi. Nessun altro
  Doria lo sa.

  \medskip\noindent\textcolor{SecretRed}{\textbf{Agenda:}}
  \textit{Completare il Colpo Riservato a qualsiasi costo personale. Se
  necessario, sacrificare la propria reputazione o un compagno — ma non le
  navi. Mai le navi.}
\end{secretbox}

\clearpage

% ════════════════════════════════════════════════════════════
% PG 2 — SER FULVIO DORIA
% @rpg-schema: <https://rpg-schema.org/instances/genova1507/pg-fulvio-doria>
%   a fitd:Character ;
%   rdfs:label "Ser Fulvio Doria" ;
%   fitd:crew <.../crew-doria> ;
%   fitd:role "Il Luogotenente / Il Braccio Destro" ;
%   fitd:action [ fitd:actionName "Survey" ; fitd:dots 3 ] ;
%   fitd:action [ fitd:actionName "Prowl" ; fitd:dots 3 ] ;
%   fitd:action [ fitd:actionName "Sway" ; fitd:dots 3 ] .
% ════════════════════════════════════════════════════════════

\pgheader[DoriaBlue]{Ser Fulvio Doria}{Il Luogotenente / Il Braccio Destro}{Doria}
\pgquote{Un uomo che non mente non sopravvive in questa città. Io sopravvivo.}

\section*{Background \& Motivazione}

\begin{rulebox}{}
  \textbf{Background:} Quarantacinque anni, cugino di secondo grado del capofamiglia.
  Ha gestito le operazioni sporche della famiglia per vent'anni — i pagamenti ai
  doganieri, le sparizioni scomode, le trattative non verbalizzabili. Sa dove sono
  sepolti i segreti di Genova. Letteralmente.

  \medskip\noindent\textcolor{DoriaBlue}{\textbf{Motivazione:}} Fedele ai Doria, ma con
  riserva mentale: la fedeltà ha un prezzo, e qualcuno gli deve ancora molte cose.
  La rivolta è un'opportunità per riscuoterle.
\end{rulebox}

\section*{Attributi \& Azioni}

\begin{attributoblock}[DoriaBlue]{INSIGHT — Mente}
  \azione{Caccia}{Hunt}{2}{Trova persone, segue tracce}
  \azione{Studio}{Study}{1}{Legge documenti, valuta testimonianze}
  \azione{Osserva}{Survey}{3}{Non gli sfugge nulla — soprattutto le menzogne}
  \azione{Ingegno}{Tinker}{1}{Scassinatore occasionale}
\end{attributoblock}
\begin{attributoblock}[DoriaBlue]{PROWESS — Corpo}
  \azione{Finezza}{Finesse}{2}{Borseggio, veleno, mani veloci}
  \azione{Ombra}{Prowl}{3}{Sorveglianza, pedinamento, sparire in una folla}
  \azione{Combattimento}{Skirmish}{2}{Coltello, scontri ravvicinati sporchi}
  \azione{Forza bruta}{Wreck}{0}{Non è la sua specialità}
\end{attributoblock}

\begin{attributoblock}[DoriaBlue]{RESOLVE — Spirito}
    \azione{Percezione}{Attune}{1}{Sente le tensioni in una stanza}
  \azione{Comando}{Command}{1}{Impartisce ordini ai subalterni}
  \azione{Relazioni}{Consort}{2}{Taverne, bordelli, i posti dove si parla}
  \azione{Persuasione}{Sway}{3}{Convince, minaccia, ricatta con eleganza}
  \end{attributoblock}

\medskip
\begin{pgclockbox}{Stress \hfill \clock{9}{0}}
  \small\textit{Traumi: $\square$ Freddo \quad $\square$ Duro \quad $\square$ Ossessionato \quad $\square$ Paranoico \quad $\square$ Irrazionale \quad $\square$ Spezzato}
\end{pgclockbox}

\section*{Abilità Speciale}
\begin{doriabox}{L'Informatore nell'Ombra}
  Una volta per sessione, Fulvio può dichiarare di aver già predisposto una
  risorsa utile nella scena attuale — un testimone corrotto, una via di fuga,
  un documento compromettente — senza che fosse stato annunciato prima.

  \medskip\noindent\textit{Nessun tiro necessario. La risorsa esiste. Il GM determina come entra in gioco.}
\end{doriabox}

\section*{Legami}
\begin{table}[h]
\centering\renewcommand{\arraystretch}{1.4}
\begin{tabularx}{\linewidth}{@{}L{2.4cm} X@{}}
  \toprule\rowcolor{DoriaBlue}
  \textcolor{white}{\textbf{PG}} & \textcolor{white}{\textbf{Legame}} \\
  \midrule
  Ottaviano & Mi crede fedele oltre ogni dubbio. Ha quasi ragione. C'è una cosa che non sa. \\
  \rowcolor{DoriaBlue!8}
  Bianca & La nipote mi spaventa. Vede troppo. Ha già capito qualcosa. \\
  Fra' Matteo & Lo uso per i messaggi più delicati. Sa che se parlasse, non rivedrebbe il suo convento. \\
  \rowcolor{DoriaBlue!8}
  Giacomo & Siamo cugini veri, di sangue. Lo proteggerei a dispetto di qualsiasi ordine. \\
  \bottomrule
\end{tabularx}
\end{table}

\section*{Equipaggiamento}
\begin{goldlist}
  \item Stiletto nascosto nello stivale destro
  \item Tre documenti falsi con identità diverse, pronti all'uso
  \item Lista criptata di debitori e creditori della famiglia — solo lui sa leggerla
  \item Chiave di una casa nel sestiere di Soziglia che ufficialmente non esiste
\end{goldlist}

\clearpage

\begin{secretbox}
  \textbf{Segreto:} Fulvio è in contatto con gli Spinola. Non li serve — li usa. Li ha
  convinti di essere un informatore quando in realtà filtra le informazioni che vuole
  far arrivare loro. Ma il gioco si è complicato: l'ultima informazione che ha passato
  era vera e ha danneggiato i Doria.

  \medskip\noindent\textcolor{SecretRed}{\textbf{Agenda:}}
  \textit{Garantirsi una via di uscita personale qualunque sia l'esito della rivolta.
  Se i Doria vincono, bene. Se perdono, Fulvio ha già trattato la propria
  sopravvivenza con una terza parte.}
\end{secretbox}

\clearpage

% ════════════════════════════════════════════════════════════
% PG 3 — BIANCA DORIA
% @rpg-schema: <https://rpg-schema.org/instances/genova1507/pg-bianca-doria>
%   a fitd:Character ;
%   rdfs:label "Bianca Doria" ;
%   fitd:crew <.../crew-doria> ;
%   fitd:role "La Stratega / L'Erede Ombra" ;
%   fitd:action [ fitd:actionName "Study" ; fitd:dots 4 ] ;
%   fitd:action [ fitd:actionName "Survey" ; fitd:dots 3 ] ;
%   fitd:action [ fitd:actionName "Consort" ; fitd:dots 3 ] ;
%   fitd:action [ fitd:actionName "Sway" ; fitd:dots 3 ] .
% ════════════════════════════════════════════════════════════

\pgheader[DoriaBlue]{Bianca Doria}{La Stratega / L'Erede Ombra}{Doria}
\pgquote{Mio nonno comanda le galee. Io comando mio nonno.}

\section*{Background \& Motivazione}
\begin{rulebox}{}
  \textbf{Background:} Ventisette anni, nipote di Ottaviano, figlia del fratello maggiore
  morto in mare. Non avrebbe dovuto contare nulla — è una donna in un mondo di uomini.
  Invece conta tutto: è lei che ha pianificato il Colpo Riservato, è lei che ha
  identificato la rotta del convoglio francese. Ottaviano la chiama «la mia nipotina»
  davanti agli altri.

  \medskip\noindent\textcolor{DoriaBlue}{\textbf{Motivazione:}} Dimostrare che il casato
  Doria può prosperare anche quando comandano le persone giuste, non solo quelle con la
  spada più lunga.
\end{rulebox}

\section*{Attributi \& Azioni}
\begin{attributoblock}[DoriaBlue]{INSIGHT — Mente}
  \azione{Caccia}{Hunt}{1}{Cerca informazioni metodicamente}
  \azione{Studio}{Study}{4}{Analisi, lingue (latino, francese, arabo), cartografia}
  \azione{Osserva}{Survey}{3}{Legge stanze, persone, situazioni politiche}
  \azione{Ingegno}{Tinker}{2}{Codici, trappole logistiche}
\end{attributoblock}
\begin{attributoblock}[DoriaBlue]{PROWESS — Corpo}
  \azione{Finezza}{Finesse}{2}{Scrive, disegna, usa veleni se necessario}
  \azione{Ombra}{Prowl}{1}{Si muove nei salotti senza farsi notare}
  \azione{Combattimento}{Skirmish}{0}{Evita lo scontro diretto}
  \azione{Forza bruta}{Wreck}{0}{Non è la sua specialità}
\end{attributoblock}
\begin{attributoblock}[DoriaBlue]{RESOLVE — Spirito}
    \azione{Percezione}{Attune}{2}{Legge le emozioni, sente le dinamiche nascoste}
  \azione{Comando}{Command}{2}{Guida con la logica, non con il volume}
  \azione{Relazioni}{Consort}{3}{Eccelle nei salotti, nelle trattative}
  \azione{Persuasione}{Sway}{3}{Convince razionalmente o emotivamente a scelta}
  \end{attributoblock}

\medskip
\begin{pgclockbox}{Stress \hfill \clock{9}{0}}
  \small\textit{Traumi: $\square$ Freddo \quad $\square$ Duro \quad $\square$ Ossessionato \quad $\square$ Paranoico \quad $\square$ Irrazionale \quad $\square$ Spezzato}
\end{pgclockbox}

\section*{Abilità Speciale}
\begin{doriabox}{Due Mosse Avanti}
  Quando Bianca pianifica un'azione (non la esegue direttamente), può dichiarare
  una \textbf{Contingenza}: se l'azione fallisce, la contingenza si attiva
  automaticamente senza tiro.

  \medskip\noindent\textit{Va dichiarata prima del tiro dell'altro PG.
  La contingenza deve essere narrativamente plausibile. Una volta per colpo.}
\end{doriabox}

\section*{Legami}
\begin{table}[h]
\centering\renewcommand{\arraystretch}{1.4}
\begin{tabularx}{\linewidth}{@{}L{2.4cm} X@{}}
  \toprule\rowcolor{DoriaBlue}
  \textcolor{white}{\textbf{PG}} & \textcolor{white}{\textbf{Legame}} \\
  \midrule
  Ottaviano & Lo amo. Sta invecchiando. Alcune decisioni le prendo io ormai — lo sa e non lo dice. \\
  \rowcolor{DoriaBlue!8}
  Ser Fulvio & Gli voglio bene come a uno zio. Anche per questo mi spaventa: so che mente. \\
  Fra' Matteo & L'unico che non mi sottovaluta. Abbastanza vecchio da non importargli di avere torto. \\
  \rowcolor{DoriaBlue!8}
  Giacomo & Utile. Pericoloso. Non mi fido di lui. Lo uso con cautela. \\
  \bottomrule
\end{tabularx}
\end{table}

\section*{Equipaggiamento}
\begin{goldlist}
  \item Taccuino cifrato con la pianificazione completa del Colpo Riservato
  \item Lettera firmata da un senatore veneziano — favore pendente da riscuotere
  \item Abiti da uomo per muoversi nei quartieri dove le donne non entrano
  \item Piccola fiala di sonnifero (non letale) — «per le emergenze»
\end{goldlist}

\clearpage

\begin{secretbox}
  \textbf{Segreto:} Bianca ha già stabilito un contatto con i Grimaldi. Non per tradire i
  Doria — per assicurare una via di fuga alla famiglia se la rivolta fallisce. L'accordo
  prevede rifugio a Monaco. Il prezzo non lo conosce ancora nessuno — i Grimaldi aspettano
  la risposta.

  \medskip\noindent\textcolor{SecretRed}{\textbf{Agenda:}}
  \textit{Garantire la sopravvivenza del casato a lungo termine. Se il Colpo Riservato
  mette in pericolo più persone di quante ne salvi, Bianca proporrà un piano alternativo
  — anche contro il volere di Ottaviano.}
\end{secretbox}

\clearpage

% ════════════════════════════════════════════════════════════
% PG 4 — FRA' MATTEO DA RECCO
% @rpg-schema: <https://rpg-schema.org/instances/genova1507/pg-fra-matteo>
%   a fitd:Character ;
%   rdfs:label "Fra' Matteo da Recco" ;
%   fitd:crew <.../crew-doria> ;
%   fitd:role "Il Cappellano / La Coscienza" .
% ════════════════════════════════════════════════════════════

\pgheader[DoriaBlue]{Fra' Matteo da Recco}{Il Cappellano / La Coscienza}{Doria}
\pgquote{Dio perdona. Io tengo il conto.}

\section*{Background \& Motivazione}
\begin{rulebox}{}
  \textbf{Background:} Sessantun anni, domenicano, cappellano delle galee Doria da
  trent'anni. Ha visto abbastanza battaglie navali, epidemie e confessioni di moribondi
  da non farsi più illusioni sulla natura umana. Crede ancora in Dio — non è sicuro
  di credere nell'umanità.

  \medskip\noindent\textcolor{DoriaBlue}{\textbf{Motivazione:}} Un popolo cristiano
  non può essere governato da potenze straniere senza il consenso della Chiesa. Il Papa
  non ha benedetto la dominazione francese. Quindi la rivolta è giusta. La logica è
  ineccepibile.
\end{rulebox}

\section*{Attributi \& Azioni}
\begin{attributoblock}[DoriaBlue]{INSIGHT — Mente}
  \azione{Caccia}{Hunt}{0}{Non è il suo campo}
  \azione{Studio}{Study}{3}{Teologia, medicina, latino e greco}
  \azione{Osserva}{Survey}{2}{Osserva le persone come leggesse le loro anime}
  \azione{Ingegno}{Tinker}{1}{Preparazione di medicamenti}
\end{attributoblock}
\begin{attributoblock}[DoriaBlue]{PROWESS — Corpo}
  \azione{Finezza}{Finesse}{1}{Mani ferme da chirurgo di bordo}
  \azione{Ombra}{Prowl}{0}{Non è la sua specialità}
  \azione{Combattimento}{Skirmish}{0}{Non combatte. Non più.}
  \azione{Forza bruta}{Wreck}{0}{No.}
\end{attributoblock}
\begin{attributoblock}[DoriaBlue]{RESOLVE — Spirito}
    \azione{Percezione}{Attune}{3}{Sente le dinamiche spirituali di un gruppo}
  \azione{Comando}{Command}{2}{L'autorità morale può valere più di quella militare}
  \azione{Relazioni}{Consort}{3}{Confraternite, bassi ranghi del clero}
  \azione{Persuasione}{Sway}{2}{La voce da predicatore. Usata raramente.}
  \end{attributoblock}

\medskip
\begin{pgclockbox}{Stress \hfill \clock{9}{0}}
  \small\textit{Traumi: $\square$ Freddo \quad $\square$ Duro \quad $\square$ Ossessionato \quad $\square$ Paranoico \quad $\square$ Irrazionale \quad $\square$ Spezzato}
\end{pgclockbox}

\section*{Abilità Speciale}
\begin{doriabox}{Confessione \& Assoluzione}
  Quando un PG del tavolo sta per subire Trauma, Fra' Matteo può intervenire.
  Il Trauma viene ridotto a un semplice punto di Stress — ma il PG deve
  confessare a Fra' Matteo (in scena, brevemente) cosa lo ha portato a quel punto.

  \medskip\noindent\textit{Una volta per PG per sessione. Fra' Matteo deve essere fisicamente presente.}
\end{doriabox}

\section*{Legami}
\begin{table}[h]
\centering\renewcommand{\arraystretch}{1.4}
\begin{tabularx}{\linewidth}{@{}L{2.4cm} X@{}}
  \toprule\rowcolor{DoriaBlue}
  \textcolor{white}{\textbf{PG}} & \textcolor{white}{\textbf{Legame}} \\
  \midrule
  Ottaviano & Lo conosco da quando aveva i capelli neri. Cerco di alleggerirgli il fardello. \\
  \rowcolor{DoriaBlue!8}
  Ser Fulvio & Non mi fido di lui. Si confessa dei peccati minori — mi chiedo cosa nasconda dietro. \\
  Bianca & Lei è il futuro di questa famiglia. Prego per lei ogni mattina. \\
  \rowcolor{DoriaBlue!8}
  Giacomo & Un peccatore allegro. Li preferisco ai santi tristi. Gli voglio bene come a un figlio. \\
  \bottomrule
\end{tabularx}
\end{table}

\section*{Equipaggiamento}
\begin{goldlist}
  \item Breviario con fogli sciolti nascosti dentro — messaggi cifrati
  \item Kit medico da campo: bende, alcol, ago e filo robusto
  \item Lettera del vescovo di Genova: libertà di movimento in tutti i sestieri
  \item Rosario di legno di ulivo — l'unica cosa rimasta del suo convento di origine
\end{goldlist}

\clearpage

\begin{secretbox}
  \textbf{Segreto:} Fra' Matteo sa del debito segreto di Ottaviano sulle galee. Lo ha
  saputo durante una confessione — non di Ottaviano, ma del banchiere. Il sigillo
  confessionale lo paralizza. Non può dirlo. Ma potrebbe agire indirettamente per
  risolvere il problema.

  \medskip\noindent\textcolor{SecretRed}{\textbf{Agenda:}}
  \textit{Proteggere Ottaviano da se stesso. Se il vecchio ammiraglio prende una
  decisione che lo distruggerà — sacrificarsi, consegnarsi ai francesi, qualcosa di
  eroico e inutile — Fra' Matteo farà di tutto per impedirlo, anche contro gli ordini
  espliciti.}
\end{secretbox}

\clearpage

% ════════════════════════════════════════════════════════════
% PG 5 — GIACOMO DETTO IL MORO
% @rpg-schema: <https://rpg-schema.org/instances/genova1507/pg-giacomo-moro>
%   a fitd:Character ;
%   rdfs:label "Giacomo detto il Moro" ;
%   fitd:crew <.../crew-doria> ;
%   fitd:role "Il Contrabbandiere / L'Uomo di Strada" .
% ════════════════════════════════════════════════════════════

\pgheader[DoriaBlue]{Giacomo detto il Moro}{Il Contrabbandiere / L'Uomo di Strada}{Doria}
\pgquote{Non sono dei vostri. Ma per stasera, forse sì.}

\section*{Background \& Motivazione}
\begin{rulebox}{}
  \textbf{Background:} Trentasei anni, origini oscure (madre genovese, padre saraceno —
  non smentisce perché gli torna utile). Contrabbandiere, sensale, informatore per
  chiunque paghi. I Doria lo usano per i rifornimenti non ufficiali e per le
  informazioni sui movimenti portuali francesi. Non è un membro della famiglia —
  è uno strumento con la propria agenda.

  \medskip\noindent\textcolor{DoriaBlue}{\textbf{Motivazione:}} Denaro, prima di tutto.
  Ma anche — lo ammette solo ubriaco — una Genova libera sarebbe meglio per la gente
  come lui: quella senza casato, senza nome, senza garanzie.
\end{rulebox}

\section*{Attributi \& Azioni}
\begin{attributoblock}[DoriaBlue]{INSIGHT — Mente}
  \azione{Caccia}{Hunt}{3}{Trova persone, merci, info in qualsiasi porto}
  \azione{Studio}{Study}{1}{Lingue (arabo, catalano, greco storto)}
  \azione{Osserva}{Survey}{3}{Sopravvissuto trent'anni grazie a questo}
  \azione{Ingegno}{Tinker}{2}{Barche, reti, nodi, meccanismi semplici}
\end{attributoblock}
\begin{attributoblock}[DoriaBlue]{PROWESS — Corpo}
  \azione{Finezza}{Finesse}{2}{Borseggio, scalata, movimenti precisi}
  \azione{Ombra}{Prowl}{3}{Caruggi, tetti, canali — ogni via non ufficiale}
  \azione{Combattimento}{Skirmish}{2}{Coltello, bastone, combattimento sporco}
  \azione{Forza bruta}{Wreck}{1}{Sfonda quando necessario}
\end{attributoblock}
\begin{attributoblock}[DoriaBlue]{RESOLVE — Spirito}
    \azione{Percezione}{Attune}{1}{Istinto animale per il pericolo}
  \azione{Comando}{Command}{0}{Non comanda. Non è nel suo vocabolario.}
  \azione{Relazioni}{Consort}{2}{Taverne, mercati neri, basse compagnie}
  \azione{Persuasione}{Sway}{2}{Mercanteggia, convince, bluffa}
  \end{attributoblock}

\medskip
\begin{pgclockbox}{Stress \hfill \clock{9}{0}}
  \small\textit{Traumi: $\square$ Freddo \quad $\square$ Duro \quad $\square$ Ossessionato \quad $\square$ Paranoico \quad $\square$ Irrazionale \quad $\square$ Spezzato}
\end{pgclockbox}

\section*{Abilità Speciale}
\begin{doriabox}{I Caruggi Sono Miei}
  Nel sestiere di Prè, nel porto vecchio, e nei caruggi del centro, Giacomo
  \textbf{non può essere sorpreso}. Ha sempre una via di fuga disponibile.
  Se inseguito, può far sparire sé stesso e un'altra persona senza tiro.

  \medskip\noindent\textit{Sempre attiva nei quartieri indicati. Fuori da questi,
  richiede un Tiro su Ombra.}
\end{doriabox}

\section*{Legami}
\begin{table}[h]
\centering\renewcommand{\arraystretch}{1.4}
\begin{tabularx}{\linewidth}{@{}L{2.4cm} X@{}}
  \toprule\rowcolor{DoriaBlue}
  \textcolor{white}{\textbf{PG}} & \textcolor{white}{\textbf{Legame}} \\
  \midrule
  Ottaviano & Mi rispetta come si rispetta uno strumento utile. Va bene così. \\
  \rowcolor{DoriaBlue!8}
  Ser Fulvio & Siamo cugini veri, di sangue. Lui lo dimentica quando fa comodo. Ricordo io. \\
  Bianca & L'unica Doria che mi parla come a una persona. Non so se mi piaccia o mi spaventi. \\
  \rowcolor{DoriaBlue!8}
  Fra' Matteo & Mi confesso da lui perché non si sorprende di nulla. E non parla. \\
  \bottomrule
\end{tabularx}
\end{table}

\section*{Equipaggiamento}
\begin{goldlist}
  \item Barca a remi nascosta sotto il molo di Sampierdarena — nessuno sa dove tranne lui
  \item Mappa mentale di tutte le vie sotterranee sotto il porto (non scritta da nessuna parte)
  \item Coltello da caccia con manico in osso di balena
  \item Piccolo sacchetto con monete di cinque valute diverse — sempre pronto a partire
\end{goldlist}

\clearpage

\begin{secretbox}
  \textbf{Segreto:} Giacomo ha venduto informazioni sui Doria agli Spinola. Non le più
  pericolose, ma una riguardava la rotta del convoglio francese — quella centrale nel
  Colpo Riservato. Se gli Spinola la usano, il Colpo potrebbe essere compromesso.

  \medskip\noindent\textcolor{SecretRed}{\textbf{Agenda:}}
  \textit{Sopravvivere e uscire da Genova con abbastanza denaro da ricominciare altrove.
  L'unica cosa che potrebbe fermarlo è Fulvio — se il cugino è nei guai, Giacomo resta.}
\end{secretbox}
